

\chapter{Introducción}

\section{Propósito del documento}

El presente documento contiene la documentación técnica del sistema \textbf{FisioManager}, una plataforma digital orientada a la rehabilitación física que integra terapeutas y pacientes mediante rutinas personalizadas, seguimiento continuo y retroalimentación inteligente basada en voz y con integración de Inteligencia Artificial.

El documento está dirigido a equipos de desarrollo, lideres técnicos y stakeholders del proyecto, y su objetivo es establecer un enfoque del diseño, las responsabilidades de los componentes y los requisitos del sistema.

\section{Alcance del sistema}

FisioManager es una aplicación web responsiva (con soporte móvil) desarrollada sobre tecnologías modernas. El sistema está concebido para operar en el contexto de \textbf{una única clínica de fisioterapia} (no es una plataforma multiclínica) y contempla \textbf{un único rol de Administrador} que concentra tanto la administración operativa de la clínica como la administración técnica de la plataforma. El sistema cubre los siguientes requisitos funcionales:

\begin{itemize}
  \item Gestión de usuarios (pacientes, terapeutas y el administrador único de la clínica).
  \item Gestión del historial clínico y diagnósticos.
  \item Planes de tratamiento por fases y notas clínicas de evolución.
  \item Agenda de citas con gestión de salas, horarios y recordatorios.
  \item Catálogo de ejercicios y creación de rutinas personalizadas.
  \item Asignación de rutinas y seguimiento de sesiones de ejecución.
  \item Captura de feedback del paciente mediante audio, escala de dolor y estado emocional.
  \item Transcripción y análisis con IA del feedback recibido.
  \item Generación de reportes de progreso y dashboards para terapeutas.
  \item Módulos de bienestar (meditación, respiración, contenido motivacional).
\end{itemize}

\section{Visión del producto}

FisioManager mejora el manejo y seguimiento del tratamiento, facilita la comunicación entre paciente y terapeuta, y optimiza la toma de decisiones clínicas mediante el aprovechamiento de inteligencia artificial sobre el feedback de voz del paciente. La plataforma funge como una herramienta intermedia entre las consultas presenciales, permitiendo continuidad y trazabilidad del tratamiento.

\section{Definiciones, acrónimos y abreviaturas}

\begin{longtable}{>{\bfseries\color{fmAccent}}L{2.5cm} L{12cm}}
  \toprule
  \rowcolor{fmTableHeader}
  \textbf{\color{fmAccent}Termino} & \textbf{\color{fmAccent}Definicion} \\
  \midrule
  \endhead

  API & Application Programming Interface, interfaz de programación de aplicaciones. \\
  \rowcolor{fmTableAltRow}
  BFF & Backend For Frontend, capa intermedia que adapta APIs al consumidor frontend. \\
  C4 Model & Modelo de documentación arquitectónica con cuatro niveles: Contexto, Contenedores, Componentes y Código. \\
  \rowcolor{fmTableAltRow}
  CIE-10 & Clasificación Internacional de Enfermedades, décima revisión. \\
  CRC Card & Class--Responsibilities--Collaborators. Tarjeta para diseño orientado a objetos. \\
  \rowcolor{fmTableAltRow}
  EF Core & Entity Framework Core, ORM oficial de .NET. \\
  FCM & Firebase Cloud Messaging, servicio de notificaciones push. \\
  \rowcolor{fmTableAltRow}
  IA & Inteligencia Artificial. \\
  JWT & JSON Web Token, estándar de tokens firmados para autenticación. \\
  \rowcolor{fmTableAltRow}
  LLM & Large Language Model, modelo de lenguaje de gran escala. \\
  PTT & Push To Talk, modalidad de comunicación por voz. \\
  \rowcolor{fmTableAltRow}
  RNF & Requerimiento No Funcional. \\
  RF & Requerimiento Funcional. \\
  \rowcolor{fmTableAltRow}
  SPA & Single Page Application, aplicación web de página única. \\
  STT & Speech-To-Text, transcripción de voz a texto. \\
  \rowcolor{fmTableAltRow}
  TTS & Text-To-Speech, síntesis de voz a partir de texto. \\
  UML & Unified Modeling Language, lenguaje unificado de modelado. \\
  \rowcolor{fmTableAltRow}
  CRR & Component Responsibility Report, reporte de responsabilidades de componentes. \\
  \bottomrule
\end{longtable}

\section{Uso de herramientas de Inteligencia Artificial}

Durante el desarrollo del proyecto se emplearon herramientas de Inteligencia Artificial (\textbf{Claude} y \textbf{Gemini}) como apoyo puntual en tareas de carácter técnico, con el fin de eficientar el tiempo del equipo. En conjunto, su uso representó aproximadamente el \textbf{20\%} del esfuerzo total del desarrollo del proyecto y se limitó a tres actividades concretas:

\begin{itemize}
  \item \textbf{Estructuración del directorio de la documentación}, organizando los artefactos por fase del proceso de análisis y diseño.
  \item \textbf{Dockerización del proyecto}, para facilitar la configuración y el despliegue del entorno de desarrollo.
  \item \textbf{Resolución de errores (bugs)} detectados durante la implementación.
\end{itemize}

\noindent El análisis, las decisiones de diseño y el contenido técnico de este documento son autoría del equipo.

\subsection{Inteligencia Artificial integrada en el producto}

Además del apoyo durante el desarrollo, el propio sistema FisioManager incorpora componentes de Inteligencia Artificial para el procesamiento del feedback de voz del paciente. A diferencia de las herramientas anteriores, estas se ejecutan \textbf{de forma local}, sin depender de proveedores externos en la nube, lo que favorece la privacidad de los datos clínicos:

\begin{itemize}
  \item \textbf{Whisper} (modelo de transcripción de voz, ejecutado en Python dentro del propio sistema): convierte automáticamente a texto en español el audio que graba el paciente al registrar su feedback (RF-040).
  \item \textbf{Ollama} (LLM local): genera el resumen clínico y el análisis de sentimiento del feedback transcrito, facilitando su lectura por parte del terapeuta (RF-041, RF-044).
\end{itemize}

\noindent Ambos resultados (transcripción y resumen) se almacenan asociados al feedback del paciente. La capa de IA del sistema procesa el audio de forma asíncrona, de modo que el paciente no espera la respuesta del modelo.

\subsection{Ejemplos de prompts}

A modo de evidencia del tipo de interacción que se tuvo con las herramientas de IA durante el desarrollo, se incluyen algunos ejemplos representativos de los prompts empleados:

\begin{infobox}[title=Resolución de errores (bugs)]
\textit{``Este endpoint REST devuelve un error 500 al subir el audio del feedback del paciente. Te comparto el controlador y el log. Ayúdame a identificar qué está fallando en el manejo del archivo.''}
\end{infobox}

\begin{infobox}[title=Dockerización del proyecto]
\textit{``Tengo una aplicación web con frontend, un BFF y una base de datos. Quiero dockerizarla para estandarizar el entorno de desarrollo del equipo. Genérame un Dockerfile para cada servicio y un docker-compose.yml que los maneje, incluyendo la base de datos con volumen persistente.''}
\end{infobox}

\begin{infobox}[title=Dockerización del proyecto]
\textit{``Explícame cómo configurar variables de entorno y redes en docker-compose para que el frontend se comunique con el BFF y este con la base de datos, sin exponer puertos innecesarios al host.''}
\end{infobox}
